\subsection{嘉当子代数与正则元}

定理 1. 设 \(\mathfrak{g}\) 是一个李代数。

\begin{enumerate}
\def\labelenumi{(\roman{enumi})}
\item
  若 \(x\) 是 \(\mathfrak{g}\) 的正则元，则 \(\mathfrak{g}^0(x)\) 是 \(\mathfrak{g}\) 的嘉当子代数。
\item
  若 \(\mathfrak{h}\) 是 \(\mathfrak{g}\) 的一个极大幂零子代数，且 \(x \in \mathfrak{h}\) 在 \(\mathfrak{g}\) 中正则，则 \(\mathfrak{h} = \mathfrak{g}^0(x)\) 。
\item
  若 \(\mathfrak{h}\) 是 \(\mathfrak{g}\) 的嘉当子代数，则 \(\dim (\mathfrak{h})\geq \mathrm{rk}(\mathfrak{g})\)
\item
  \(\mathfrak{g}\) 的维数为 \(\mathrm{rk}(\mathfrak{g})\) 的嘉当子代数恰是诸 \(\mathfrak{g}^0 (x)\) ，其中 \(x\) 是正则元。
\end{enumerate}

设 \(x\) 是 \(\mathfrak{g}\) 的正则元，并设 \(\mathfrak{h} = \mathfrak{g}^0(x)\) 。显然 \(\mathfrak{h}^0(x) = \mathfrak{h}\) 。由于 \(x\) 在 \(\mathfrak{h}\) 中正则（命题 9），\(\operatorname{rk}(\mathfrak{h}) = \dim(\mathfrak{h})\) ，故 \(\mathfrak{h}\) 幂零。另一方面，\(\mathfrak{h} = \mathfrak{g}^0(x) \supset \mathfrak{g}^0(\mathfrak{h}) \supset \mathfrak{h}\) ，故 \(\mathfrak{h} = \mathfrak{g}^0(\mathfrak{h})\) 是 \(\mathfrak{g}\) 的嘉当子代数（命题 4）。这就证明了 (i)。

若 \(\mathfrak{h}\) 是 \(\mathfrak{g}\) 的一个极大幂零子代数，且 \(x\in \mathfrak{h}\) 在 \(\mathfrak{g}\) 中正则，则 \(\mathfrak{h}\subset \mathfrak{g}^0 (x)\) ，且由 (i) \(\mathfrak{g}^0 (x)\) 幂零，故 \(\mathfrak{h} = \mathfrak{g}^0 (x)\) ，这就证明了 (ii)。

若 \(\mathfrak{h}\) 是 \(\mathfrak{g}\) 的嘉当子代数，则存在 \(x \in \mathfrak{h}\) 使 \(\mathfrak{h} = \mathfrak{g}^0(x)\) （命题 4 的推论 1），故 \(\dim(\mathfrak{h}) \geq \operatorname{rk}(\mathfrak{g})\) ，这就证明了 (iii)。若此外 \(\dim(\mathfrak{h}) = \operatorname{rk}(\mathfrak{g})\) ，则 \(x\) 正则。最后，若 \(x'\) 在 \(\mathfrak{g}\) 中正则，则由 (i) \(\mathfrak{g}^0(x')\) 是嘉当子代数，且显然维数为 \(\operatorname{rk}(\mathfrak{g})\) 。这就证明了 (iv)。

我们将在 §3 定理 2 中看到，当 \(k\) 的特征数为零时，\(\mathfrak{g}\) 的所有嘉当子代数都有维数 \(\mathrm{rk}(\mathfrak{g})\) 。

推论 1. 每个李代数 g 都有嘉当子代数，且 g 的秩是嘉当子代数的最小维数。

推论 2. 设 \(f: \mathfrak{g} \to \mathfrak{g}'\) 是李代数的一个满射同态。若 \(\mathfrak{h}'\) 是 \(\mathfrak{g}'\) 的嘉当子代数，则存在一个嘉当子代数 \(\mathfrak{h}\) （在 \(\mathfrak{g}\) 中）使 \(\mathfrak{h}' = f(\mathfrak{h})\) 。

设 \(\mathfrak{a} = f^{-1}(\mathfrak{h}')\) 。由推论 1，\(\mathfrak{a}\) 有一个嘉当子代数 \(\mathfrak{h}\) 。由命题 4 的推论 2，\(f(\mathfrak{h}) = \mathfrak{h}'\) 。我们证明 \(\mathfrak{h}\) 是 \(\mathfrak{g}\) 的嘉当子代数。设 \(\mathfrak{n}\) 是 \(\mathfrak{h}\) 在 \(\mathfrak{g}\) 中的正规化子。只需证明 \(\mathfrak{h} = \mathfrak{n}\) 。若 \(x \in \mathfrak{n}\) ，则 \(f(x)\) 属于 \(\mathfrak{h}'\) 在 \(\mathfrak{g}'\) 中的正规化子，故 \(f(x) \in \mathfrak{h}'\) 且 \(x \in \mathfrak{a}\) ；但 \(\mathfrak{h}\) 在 \(\mathfrak{a}\) 中是自身的正规化子，故 \(x \in \mathfrak{h}\) 。

推论 3. 每个李代数 \(\mathfrak{g}\) 都是它的嘉当子代数之和。

g 的诸嘉当子代数之和 s 包含 g 的正则元之集（定理 1 (i)）。由于该集在 g 中按 Zariski 拓扑稠密，故 s = g。

命题 10. 设 g 是一个李代数，a 是 g 的一个交换子代数，c 是 a 在 g 中的交换子。假定 \(ad_{g}x\) 对所有 \(x \in a\) 半单。那么 c 的嘉当子代数恰是 g 的包含 a 的嘉当子代数。

设 \(\mathfrak{h}\) 是 \(\mathfrak{c}\) 的嘉当子代数。由于 \(\mathfrak{a}\) 包含在中心 \(\mathfrak{z}\) 中，即有 \(\mathfrak{c}, \mathfrak{a} \subset \mathfrak{z} \subset \mathfrak{h}\) （命题 5）。设 \(\mathfrak{n}\) 是 \(\mathfrak{h}\) 在 \(\mathfrak{g}\) 中的正规化子。那么

\[
[ \mathfrak {a}, \mathfrak {n} ] \subset [ \mathfrak {h}, \mathfrak {n} ] \subset \mathfrak {h}.
\]

由于诸 \(ad_{g}x\) （\(x \in a\) ）半单且两两交换，由《代数》第八章 §5 第 1 节可知，存在 n 的一个在 \(ad_{g}a\) 下稳定的向量子空间 d 使得 \(n = h \oplus d\) 。于是 \([a, d] \subset h \cap d = 0\) ，故 \(d \subset c\) 。从而 n 是 h 在 c 中的正规化子，因此 \(n = h\) ，故 h 是 g 的包含 a 的嘉当子代数。

反过来，设 \(\mathfrak{h}\) 是 \(\mathfrak{g}\) 的包含 \(\mathfrak{a}\) 的嘉当子代数。那么 \(\mathfrak{h} = \mathfrak{g}^0 (\mathfrak{h})\subset \mathfrak{g}^0 (\mathfrak{a})\) ，且由假定 \(\mathfrak{g}_0(\mathfrak{a}) = \mathfrak{g}^0 (\mathfrak{a}) = \mathfrak{c}\) 。因此 \(\mathfrak{a}\subset \mathfrak{h}\subset \mathfrak{c}\) ，且 \(\mathfrak{h}\) 是 \(\mathfrak{c}\) 的嘉当子代数（因为它在 \(\mathfrak{g}\) 中等于自身的正规化子，故在 \(\mathfrak{c}\) 中更其如此）。

命题 11. 设 \(\mathfrak{n}\) 是李代数 \(\mathfrak{g}\) 的一个幂零子代数。存在 \(\mathfrak{g}\) 的一个包含于 \(\mathfrak{g}^0 (\mathfrak{n})\) 的嘉当子代数。

设 \(\mathfrak{a} = \mathfrak{g}^0 (\mathfrak{n})\) 。于是 \(\mathfrak{n}\subset \mathfrak{a}\) ，由于 \(\mathfrak{n}\) 幂零。若 \(x\in \mathfrak{a}\) ，设 \(\mathrm{P}(x)\) 是 \(\mathfrak{g} / \mathfrak{a}\) 的自同态（由 ad \(x\) 定义）的行列式。用 \(\mathfrak{a}'\) 记 \(x\in \mathfrak{a}\) 中使 \(\mathrm{P}(x)\neq 0\) 者所成之集，它是 \(\mathfrak{a}\) 中按 Zariski 拓扑的一个开子集；关系 \(x\in \mathfrak{a}'\) 与 \(\mathfrak{g}^0 (x)\subset \mathfrak{a}\) 等价。由 §1 第 2 节的命题 7 (ii)，存在 \(y\in \mathfrak{n}\) 使 \(\mathfrak{g}^0 (y) = \mathfrak{a}\) ，且 \(y\in \mathfrak{a}'\) ，故 \(\mathfrak{a}'\) 非空。由于 \(\mathfrak{a}'\) 开，它与 \(\mathfrak{a}\) 的正则元之集的交非空。设 \(x\) 是该交中的一个元素。那么 \(\mathfrak{g}^0 (x)\subset \mathfrak{a}\) ，且 \(\mathfrak{g}^0 (x)\) 是 \(\mathfrak{a}\) 的嘉当子代数，故幂零。另一方面，§1 第 3 节的命题 10 (i) 表明 \(\mathfrak{g}^0 (x)\) 在 \(\mathfrak{g}\) 中是自身的正规化子；因此它是 \(\mathfrak{g}\) 的嘉当子代数，证毕。

